\section{Indulgences and the Church's Suffrages for the Holy Souls}

Death alters a Christian's state, but it never breaks the Mystical Body of Christ. Pilgrims on earth (the Church Militant), souls undergoing purification (the Church Suffering), and saints beholding God (the Church Triumphant) remain one Body under one Head in the Communion of Saints. St.\ Paul declares this organic solidarity: ``And if one member suffer any thing, all the members suffer with it: or if one member glory, all the members rejoice with it'' (1 Cor 12:26); indeed, he rejoices to ``fill up those things that are wanting of the sufferings of Christ, in my flesh, for his body, which is the church'' (Col 1:24).

In \textit{Spe salvi}, Pope Benedict XVI reflects on how deeply this solidarity crosses the grave: because no one lives, sins, or is saved alone, our thoughts, words, and deeds continually flow into the lives of others for good or ill, so that loving intercession can still reach the departed and is never offered in vain.\endnote{Pope Benedict XVI, Encyclical Letter \textit{Spe salvi} (30 November 2007), \S 48.}

Because the Holy Souls can no longer merit or satisfy for themselves yet remain Christ's living members, God ordains that the charity of the living reach them with real relief. These works of intercession and satisfaction offered for the dead are called \textit{suffrages} (from \textit{suffragium}, pleading aid).

Long before scholastic treatises analyzed suffrages, the liturgy practiced them as an apostolic inheritance, continuing Israel's intercession in Job's sacrifices for his children (Job 1:5), Tobias's alms at burial (Tob 4:18), and Judas Machabeus's sin-offering for the fallen (2 Mach 12:43--46). East and West bear identical testimony from antiquity.

In fourth-century Jerusalem, St.\ Cyril describes the Eucharistic Anaphora to the newly baptized: after the consecration, over ``that sacrifice of propitiation,'' the Church commemorates the saints and then prays for the departed:

\begin{quote}
Then we commemorate also those who have fallen asleep before us, first Patriarchs, Prophets, Apostles, Martyrs, that at their prayers and intercessions God would receive our petition. Then on behalf also of the Holy Fathers and Bishops who have fallen asleep before us, and in a word of all who in past years have fallen asleep among us, believing that it will be a very great benefit to the souls, for whom the supplication is put up, while that holy and most awful sacrifice is set forth.\endnote{St.\ Cyril of Jerusalem, \textit{Catechetical Lectures} XXIII (Mystagogical V), \S 9, in \textit{NPNF} II, vol.\ 7.}
\end{quote}

To skeptics asking how remembering the dead at Mass aids them, Cyril compares us to kinsmen offering a crown to a king for banished offenders, except that we ``weave no crown, but offer up Christ sacrificed for our sins, propitiating our merciful God for them as well as for ourselves.''\endnote{St.\ Cyril of Jerusalem, \textit{Catechetical Lectures} XXIII, \S 10, in \textit{NPNF} II, vol.\ 7.}

St.\ John Chrysostom speaks with equal force in his Homilies on Philippians:

\begin{quote}
Not in vain did the Apostles order that remembrance should be made of the dead in the dreadful Mysteries. They know that great gain resulteth to them, great benefit; for when the whole people stands with uplifted hands, a priestly assembly, and that awful Sacrifice lies displayed, how shall we not prevail with God by our entreaties for them?\endnote{St.\ John Chrysostom, \textit{Homilies on Philippians} III, in \textit{NPNF} I, vol.\ 13.}
\end{quote}

Preaching on First Corinthians, he urges the faithful to aid the departed ``not by tears, but by prayers and supplications and alms and offerings,'' for ``the common Expiation of the world is even before us \dots\ For in truth one body are we all, though some members are more glorious than others; and it is possible from every source to gather pardon for them, from our prayers, from our gifts in their behalf, from those whose names are named with theirs.''\endnote{St.\ John Chrysostom, \textit{Homilies on First Corinthians} XLI.8, in \textit{NPNF} I, vol.\ 12.}

In the Latin West, St.\ Monica's dying words at Ostia in Book IX of St.\ Augustine's \textit{Confessions} stand unmatched: ``Lay this body anywhere, let not the care for it trouble you at all. This only I ask, that you will remember me at the Lord's altar, wherever you be.''\endnote{St.\ Augustine, \textit{Confessions} IX.11.27, in \textit{NPNF} I, vol.\ 1.} After ``the sacrifice of our redemption was offered up unto Thee for her'' beside her grave (IX.12.32), Augustine prayed:

\begin{quote}
And woe even unto the praiseworthy life of man, if, putting away mercy, Thou shouldest investigate it. \dots\ I then, O my Praise and my Life, Thou God of my heart, putting aside for a little her good deeds, for which I joyfully give thanks to Thee, do now beseech Thee for the sins of my mother. \dots\ Forgive her, O Lord, forgive her, I beseech Thee; ``enter not into judgment'' with her. Let Thy mercy be exalted above Thy justice \dots\ And inspire, O my Lord my God, inspire Thy servants my brethren \dots\ that so many of them as shall read these confessions may at Thy altar remember Monica, Thy handmaid, together with Patricius, her sometime husband.\endnote{St.\ Augustine, \textit{Confessions} IX.13.34--35, 37, in \textit{NPNF} I, vol.\ 1.}
\end{quote}

In \textit{The City of God}, Augustine specifies whom these suffrages aid: neither the obstinately damned nor the already perfected saints, but ``those who, having been regenerated in Christ, did not spend their life so wickedly that they can be judged unworthy of such compassion, nor so well that they can be considered to have no need of it.''\endnote{St.\ Augustine, \textit{The City of God} XXI.24, trans.\ Dods (1871).}

How can one Christian satisfy another's temporal debt? St.\ Thomas Aquinas explains that while no one bears another's personal guilt, union of charity joins wills into one: thus ``since those who differ as to the debt of punishment, may be one in will by the union of love, it happens that one who has not sinned, bears willingly the punishment for another,'' as Christ Himself satisfied for us.\endnote{Aquinas, \textit{ST} I-II, q.\ 87, a.\ 7, corpus and ad 3; a.\ 8, corpus.} The \textit{Roman Catechism} sets out this privilege plainly:

\begin{quote}
But in this the mercy and goodness of God shine conspicuous, and demand our grateful acknowledgments, that he has granted to our frailty the privilege that one may satisfy for another. This, however, is a privilege which is confined to the satisfactory part of penance alone, and extends not to contrition and confession: no man can be contrite or confess for another; whilst those who are gifted with divine grace may pay through others what is due to the divine justice, and thus we may be said in some measure to bear each other's burdens. \dots\ As then the foot does not perform its functions solely for itself, but also for sake of the other members \dots\ so works of satisfaction are common to all the members of the Church.\endnote{\textit{Catechism of the Council of Trent}, Part II, On the Sacrament of Penance (Satisfaction), trans.\ Donovan, p.\ 204.}
\end{quote}

The Catechism adds a precise safeguard: personal works of satisfaction require sanctifying grace, whereas even a sinner may beg God's mercy or arrange for a Mass, since the Mass draws its power principally from Christ the High Priest.

Following the Profession of Faith at Lyons II (DS 856) and the ecumenical definitions of Florence (DS 1304) and Trent (DS 1820), the \textit{Baltimore Catechism No.\ 3} summarizes the Church's suffrages concisely: ``The faithful on earth can help the souls in Purgatory by their prayers, fasts, almsdeeds; by indulgences, and by having Masses said for them.''\endnote{\textit{Baltimore Catechism No.\ 3} (1921), Q.\ 1385, p.\ 308.}

Foremost stands the \textit{Holy Sacrifice of the Mass}, by which the Holy Souls are assisted above all (Trent, DS 1820; cf.\ DS 1743, 1753), since Calvary's spotless Victim is offered directly to the Father. Next stand \textit{prayer}---both the Church's liturgical offices and private devotion---\textit{almsgiving and works of mercy}, and \textit{fasting, mortification, and patient suffering} united to the Cross.

Beside these stands a distinctly ecclesial treasury of mercy: \textit{Indulgences}. Rooted in the power of the keys---``whatsoever thou shalt bind upon earth, it shall be bound also in heaven: and whatsoever thou shalt loose on earth, it shall be loosed also in heaven'' (Matt 16:19; cf.\ 18:18; 2 Cor 2:10)---the early Church already shortened canonical penances through the intercession of the martyrs. At Session XXV (1563), the Council of Trent defined in its \textit{Decree on Indulgences} that Christ granted the Church power to confer indulgences, that she has exercised this power from antiquity, and that their use is deeply salutary for Christian people, while commanding bishops to abolish mercenary traffic and all disorderly abuses.\endnote{Council of Trent, Session XXV (1563), \textit{Decree on Indulgences} (DS 1835).}

In the Apostolic Constitution \textit{Indulgentiarum doctrina} (1 January 1967)---codified for the Latin Church in the 1983 \textit{Code of Canon Law} (cann.\ 992--997) and reiterated in the \textit{Catechism of the Catholic Church}---Pope St.\ Paul VI restated this doctrine and reformed the Church's practical indulgence discipline. An indulgence is the remission before God of the temporal punishment due to sins whose guilt has already been forgiven, gained by a duly disposed member of the faithful under prescribed conditions through the ministry of the Church, which as dispenser of redemption applies with authority the treasury of the satisfactions of Christ and the saints.\endnote{Pope St.\ Paul VI, Apostolic Constitution \textit{Indulgentiarum doctrina} (1 January 1967), Norm 1; 1983 \textit{Code of Canon Law}, can.\ 992; \textit{Catechism of the Catholic Church}, \S 1471.}

Four clarifications dispel every caricature and distinguish enduring doctrine from post-1967 discipline. First, an indulgence is \textit{never} pardon for sin or license to sin; as the \textit{Baltimore Catechism No.\ 3} states, ``An Indulgence is not a pardon of sin, nor a license to commit sin, and one who is in a state of mortal sin cannot gain an Indulgence.''\endnote{\textit{Baltimore Catechism No.\ 3} (1921), Q.\ 841, pp.\ 177--178.} Second, an indulgence is either \textit{partial} or \textit{plenary} (1983 \textit{Code of Canon Law}, can.\ 993; \textit{Baltimore Catechism No.\ 3}, Qq.\ 843--845, p.\ 178). Under the discipline reformed by Pope St.\ Paul VI in \textit{Indulgentiarum doctrina} (1 January 1967, Norms 4--5, 7, 20)---codified in the 1983 \textit{Code of Canon Law} (cann.\ 992--997) and implemented in the Apostolic Penitentiary's \textit{Enchiridion Indulgentiarum} (4th ed., 1999, Norm 20)---the Church abolished the older determination of partial indulgences in ``days'' and ``years'' corresponding to ancient canonical penances (as still explained in the 1921 \textit{Baltimore Catechism No.\ 3}, Q.\ 863, p.\ 183). Instead, for a partial indulgence, the Church by her intervention grants a remission equal to what the faithful already obtain by the pious work itself, effectively doubling that remission. For a plenary indulgence under current Latin-Church discipline, the required work must be joined to three usual conditions already summarized in the \textit{Baltimore Catechism No.\ 3}---sacramental confession, Eucharistic Communion, and prayer for the intentions of the Supreme Pontiff---together with complete detachment from all sin, even venial sin; where that full disposition or any condition is lacking, the indulgence remains partial.\endnote{Pope St.\ Paul VI, \textit{Indulgentiarum doctrina} (1 January 1967), Norms 4--5, 7, 20; 1983 \textit{Code of Canon Law}, cann.\ 992--997; Apostolic Penitentiary, \textit{Enchiridion Indulgentiarum} (4th ed., 1999), Norm 20; \textit{Baltimore Catechism No.\ 3} (1921), Qq.\ 843--845 (p.\ 178), Q.\ 861 (p.\ 182), and Q.\ 863 (p.\ 183).}

Third, the \textit{Treasury of the Church} is neither a material storehouse nor a rival to Calvary, but the infinite, inexhaustible worth that Christ's merits hold before God, joined to the prayers and good works of the Blessed Virgin Mary and of all the saints, who in Christ's grace achieved a superabundance of spiritual goods belonging to the whole Mystical Body.\endnote{Pope St.\ Paul VI, \textit{Indulgentiarum doctrina} (1967), \S 5; cf.\ \textit{Catechism of the Catholic Church}, \S\S 1476--1477; \textit{Baltimore Catechism No.\ 3} (1921), Qq.\ 853--854, pp.\ 180--181.}

Fourth, whereas the Church grants indulgences to the living authoritatively by way of absolution, she applies them to the departed strictly \textit{per modum suffragii}, by way of suffrage.\endnote{Pope St.\ Paul VI, \textit{Indulgentiarum doctrina} (1967), Norm 3; 1983 \textit{Code of Canon Law}, can.\ 994; \textit{Catechism of the Catholic Church}, \S\S 1471, 1479.} Because the dead lie beyond earthly jurisdiction, the Church does not command God; she presents a share from the Treasury of Christ and His saints and implores His clemency for their release.

Thus Catholic faith turns grief into active charity. Because no prayer of love is wasted---if the soul we name already beholds God, He applies our suffrage to another in need\endnote{Cf.\ \textit{Baltimore Catechism No.\ 3} (1921), Q.\ 1384, p.\ 308.}---every Mass, prayer, fast, almsdeed, and indulgence binds earth to heaven in living mercy.
